\section{Conclusion and Future Work}
\label{submission:conclusion}

In this work, we have established a mathematically rigorous foundation for dynamic model merging, addressing the systemic lack of theoretical justification in current input-dependent routing protocols. 
Through the lens of statistical learning theory, we formally analyzed the generalization capabilities of dynamic parameter-space blending. 
By deriving the empirical Rademacher complexity bounds of the layer-wise dynamic merged hypothesis class, we proved that the generalization error of the router is bounded by a weighted Frobenius norm of its weights, scaled by task-specific expert activation energies.

This learning-theoretic analysis led directly to the formulation of \textbf{R2D-Merge} and its novel regularization mechanism, \textbf{Covariance-weighted Frobenius Regularization (CFR)}. 
CFR adaptively penalizes routing parameters based on localized expert weight sensitivity and feature covariances, pre-computed offline on an extremely sparse calibration split ($N=64$) and introducing \emph{zero online computational overhead} during inference. 

Our empirical evaluations on a Vision Transformer backbone across four vision datasets (MNIST, FashionMNIST, CIFAR-10, and SVHN) validate our theoretical approach. 
While unregularized routers and complex quantum-inspired heuristics suffer catastrophic performance collapse (up to -13.00\% drops) when dynamic coefficients are averaged to satisfy single-model hardware constraints on heterogeneous streams, R2D-Merge demonstrates absolute resilience (0.00\% drop). It outperforms unregularized baselines by up to 11.50\% absolute accuracy in the collapsed state. 
Additionally, CFR outperforms uniform L2 decay on complex datasets, highlighting the advantage of task-adaptive, covariance-weighted regularization.

\textbf{Future Directions.} There are several promising avenues for future research. 
First, we plan to extend our Rademacher generalization bounds to non-linear dynamic model blending, where routing networks incorporate multi-layer perceptrons or attention mechanisms. 
Second, we aim to address the expressive scaling of the routing dimension $d$ when scaling R2D-Merge to larger architectures, such as CLIP-ViT-L or Large Language Models (LLMs). While $d=4$ or $d=8$ is highly expressive and sufficient for a ViT-Tiny, billions-of-parameter models may require larger routing dimensions (e.g., $d \geq 32$). Under such regimes, the offline-computed covariance matrices $C_{l,k}$ could become high-dimensional and suffer from singular/low-rank issues on small calibration splits. To preserve statistical and computational efficiency, we plan to explore structured covariance approximations, such as diagonal, block-diagonal, or Kronecker-factored approximate curvature (KFAC) style approximations of $C_{l, k}$. In particular, modeling $C_{l, k}$ as a Kronecker product of smaller factors would reduce estimation noise, maintain statistical tractability on extremely sparse calibration splits, and keep the offline profiling cost exceptionally low.
Third, we intend to investigate a hybrid static-dynamic merging framework, which we term \emph{TIES-Routing}. In this framework, task vectors $V_k^{(l)}$ are first pruned and sign-resolved using offline heuristics like TIES-Merging to eliminate destructive parameter-space conflicts. Subsequently, our R2D-Merge router can be optimized over the pruned task vectors. Combining sign-agreement/redundancy-reduction with dynamic, covariance-weighted routing is a highly promising direction that could unlock superior multi-task performance.
Fourth, we aim to investigate the application of R2D-Merge to autoregressive language modeling tasks, exploring how covariance-weighted regularization scales to billions of parameters in LLMs. 
Finally, we intend to study the interaction between parameter-space blending and activation-space routing to build fully unified, mathematically bounded multi-task networks.
